We begin by recalling and defining some general concepts that play an important role throughout this thesis.

Let $X$ be a Banach space and $T > 0$. A function $f: [0, T] \to X$ is called \emph{strongly measurable} if there is a sequence of simple functions $(f_n)_{n \in \N}$ that converges pointwise to $f$ at almost every $t \in [0, T]$. Now for $p \in [1, \infty]$, the space

\begin{equation*}
    L^{p}(0, T; X)
\end{equation*}

\noindent
consists of all strongly measurable functions $u$ such that

\begin{equation*}
    \norm{u}_{L^p(0, T; X)} := \left( \int_0^T \norm{u(t)}_X dt \right)^{1/p} < \infty
\end{equation*}

\noindent
for $p \in [1, \infty)$, and

\begin{equation*}
    \norm{u}_{L^\infty(0, T; X)} := \esssup_{t\ \in\ [0, T]} \norm{u(t)}_X < \infty.
\end{equation*}

\noindent
Lastly,

\begin{equation*}
    C([0, T]; X)
\end{equation*}

\noindent
denotes the space of all continuous functions $u: [0, T] \to X$ with

\begin{equation*}
    \norm{u}_{C([0, T]; X)} := \max_{t \in [0, T]} \norm{u(t)}_X < \infty.
\end{equation*}

\noindent
We can take $X$ to be, for instance, $L^2(\Omega)$. Then formally, $C([0, T]; L^2(\Omega))$ consists of continuous functions $u$ from the interval time [0, T] to $L^2(\Omega)$. For a more intuitive interpretation, you might feel inclined to interpret them as the corresponding functions $\widetilde{u}$ from $[0, T] \times \Omega$ to $\R$, given by $\widetilde{u}(t, x) = u(t)(x)$. We do need to be careful when we use this interpretation, though, as the notation might make it seem like these functions $\widetilde{u}$ are continuous in time, while this is generally not the case.

The next ideas are often used to translate local properties into global results. This method plays a fundamental part in our later discussions about the formulation and analysis of weak forms and solutions to PDEs.

Let $X$ now be a topological space. A function $\phi: X \to \R$ is said to have \emph{compact support} if the set

\begin{equation*}
    \text{supp}(\phi) := \overline{\{ x \in X \mid \phi (x) \neq 0 \}}
\end{equation*}

\noindent
is a compact subset of $X$. For any open $\Omega \subseteq \R^m$, the space of all smooth functions with compact support in $\Omega$ is denoted by $C^\infty_c(\Omega)$.

An important property of these functions is that they must be 0 on the boundary of $\Omega$. In the particular context of PDE analysis, they are often referred to as \emph{test functions}, but a more general term is \emph{variations}. They have a wide range of applications, discussed in the field of the Calculus of Variations. A large portion of its theory is built around the following result, hence its name.

\begin{lemma}[Fundamental Lemma of the Calculus of Variations]
    \label{thm:flcv}
    Let $\Omega \subseteq \R^{m}$ be open and $f \in L^{2}(\Omega)$. If

    \begin{equation*}
        \int_{\Omega} f \phi = 0
    \end{equation*}

    holds for all $\phi \in C^{\infty}_{c}(\Omega)$, then $f = 0$ almost everywhere in $\Omega$.
\end{lemma}

Multiplying with a test function and integrating is sometimes referred to as \emph{testing against} a test function $\phi \in C^\infty_c(\Omega)$. 

\begin{definition}
    When we are differentiating functions in $C^k$ defined in a subset of $\R^m$, we sometimes use the abbreviated notation

    \begin{equation*}
        D^\alpha := \frac{\partial^{|\alpha|}}{\partial x_1^{\alpha_1} \dots \partial x_m^{\alpha_m}} ,
    \end{equation*}

    where $\alpha = (\alpha_1, \dots, \alpha_m) \in \N^m$ is called a multi-index and $|\alpha| := \alpha_1 + \dots + \alpha_m \leq k$.
\end{definition}

\begin{definition}
    Let $\Omega \subseteq \R^m$ be open. A function $f: \Omega \to \R$ is said to be \emph{locally integrable} if

    \begin{equation*}
        \int_K |f| < \infty
    \end{equation*}

    for all compact subsets $K$ of $\Omega$. The space of all locally integrable functions on $\Omega$ is denoted by $L^1_\text{loc}(\Omega)$.
\end{definition}

An important fact that should be noted is that every funcion that is square-integrable is also locally integrable. That is, $L^2(\Omega) \subseteq L^{1}_\text{loc}(\Omega)$. For the sake of completeness, we will sometimes state definitions or results with the hypothesis that a function is locally integrable, only to use them on functions that are square-integrable.